\begin{frame}{Persistent Homology Transform (PHT)}
    \centering
    \animategraphics[loop,autoplay,width=.6\linewidth]{10}{topology/pht/animation/animation-}{0}{98}
\end{frame}

\frame{
    \frametitle{Defining the PHT, 2013}
\begin{block}{}
\centering
	\begin{itemize}
		\item Turner, Mukherjee, and Boyer show that a set of
                    persistence diagrams uniquely represents three-dimensional
                    shapes.\footnote{\scriptsize{Turner,
                    Mukherjee, and Boyer. Persistent Homology Transform for
                    Modeling Shapes and Surfaces. Information and Inference: A
                    Journal of the IMA, 3(4), 310-344, 2014.  Preprint:
                    arXiv:1310.1030, 2013.}}
                \item Method: locate vertices, then links.
                \item \bf{With all (infinite) directions, we have faithful
                representation}.
	\end{itemize}
\end{block}
}


\frame{
    \frametitle{PHT is Faithful}

    \begin{block}{The PHT}
        Let $K$ be a shape in $\R^d$.
        $$PHT(K) := \left\{ \left(s, D(s,K)\right) \right\}_{s \in \S^{d-1}}.$$
    \end{block}

    \pause

    \begin{block}{Theorem}
        No two distinct shapes have the same PHT (nor ECT).
    \end{block}
}


\frame{
    \frametitle{Five Years Later ...}
    \begin{block}{}
        \begin{itemize}
            \item RLB, BTF, RM, SM, DLM, DS, AS, LW. {Learning Simplicial
                Complexes from Persistence Diagrams}. CCCG, 2018.
            \item Curry, Mukherjee, and Turner. ``How
                many directions determine a shape and other sufficiency results
                for two topological transforms.'' Trans AMS B (2022): 1006--1043.
            \item Ghrist, Levanger, and Mai. ``Persistent
                homology and Euler integral transforms.'' J.\ Applied and
                Computational Topology 2.1-2 (2018): 55--60.
            \item (and a couple more recent papers)
        \end{itemize}
    \end{block}
}

\frame<1-8>[label=faithfuldisc]{
    \frametitle{Faithful Discretization}

    \begin{columns}
        \column{.45\textwidth}
            \begin{block}{Discretization}
                \onslide<2->{Let $S \subseteq \S^{d-1}$.}
                $$
                   \onslide<3->{\mathcal{S} = \left\{ \big( s,}%
                   \onslide<1->{D(s,K)}%
                   \onslide<3->{\big) \right\}_{s \in S} }%
                $$
            \end{block}

            \onslide<4->{
            \begin{block}{Faithful}
                $\mathcal{S}$ is \emph{faithful} if invertible.
            \end{block}
            }

            \onslide<5->{
            \begin{block}{Minimum Faithful Set}
                $\mu(K)$ is the size of the smallest faithful discretization
                of~$PHT(K)$.
            \end{block}
            }

        \column{.45\textwidth}
            \centering
            \only<6>{\includegraphics[height=1.5in]{topology/pht/ex-v-rep-init}}%
            \only<7>{\includegraphics[height=1.5in]{topology/pht/ex-v-rep-full}}%
            \only<8>{\includegraphics[height=1.5in]{topology/pht/ex-v-rep-pht}}%
            \only<9>{\includegraphics[height=1.5in]{topology/pht/ex-v-rep-pht-guidelines}}%
            \only<10->{\includegraphics[height=1.5in]{topology/pht/ex-v-rep-apht-guidelines}}%
    \end{columns}
}


\frame{
    \frametitle{Necessary Perpendicular Directions}

    \begin{theorem}{}
        Let $K$ be a simplicial complex in $\R^d$ and $\sigma$ a maximal
        simplex.
        \pause
        Then, a faithful set must use at least $d-\dim{\sigma}+1$ directions perpendicular to
        $\sigma$.
    \end{theorem}

}


\frame<1-4>{
    \frametitle{Tricky Elbows}

    \begin{columns}
        \column{\textwidth}
        \centering
        \begin{block}{}
            \only<1>{\includegraphics[height=2in]{topology/pht/observability-nobend}}%
            \only<2>{\includegraphics[height=2in]{topology/pht/observability-bend}}%
            \only<3>{\includegraphics[height=2in]{topology/pht/observability}}%
            \only<4>{\includegraphics[height=2in]{topology/pht/observability-flatter}}%
            \only<5->{\includegraphics[height=2in]{topology/pht/observability-tri}}%
        \end{block}

    \end{columns}
    }


\frame{
    \frametitle{Counting Directions}

    \begin{block}{}
    \begin{columns}
        \column{.5\textwidth}
            \centering
            \only<1-2>{\includegraphics[height=1.5in]{topology/pht/ex-v-perp-necc}}%
            \only<3->{\includegraphics[height=1.5in]{topology/pht/ex-v-rep-pht-guidelines}}%
        \column{.5\textwidth}
            \centering
            \only<2->{\includegraphics[height=1.5in]{topology/pht/ex-v-alt}}%
    \end{columns}
    \end{block}
}

\frame<1-2>[label=triangles]{
    \frametitle{How Many Directions do We Need?}

    \begin{block}{}
        \centering
        \includegraphics[width=\textwidth]{topology/pht/construction-edges}
        \begin{itemize}
            \pause \item $25$ persistence diagrams necessary.
            \pause \item $3$ verbose persistence diagrams necessary \& sufficient.
            \pause \item $\epsilon$-net approach: $> 1$k diagrams.
        \end{itemize}
    \end{block}
}
